\documentclass[a4paper,12pt]{article}
	\usepackage[ansinew]{inputenc}
	\usepackage{ae}
	\usepackage[brazil]{babel}
	\usepackage{lipsum}
	\usepackage{xcolor}
	\usepackage{graphicx}
	
	\graphicspath{
		{Figuras/EPS/}
		{Figuras/JPG/}
		{Figuras/PDF/}
		{Figuras/PNG/}
		}
			
\begin{document}

	%\lipsum[1-5]

	\begin{figure}[bht]
		\centering
		\includegraphics[height=0.2\textheight]{Tei-Gi}
		\caption{o Tei-Gi ou In Yang, símbolo de equanimidade.}
		\label{fig:tei-gi} % <-- Coloque o rótulo sempre após o \caption.
	\end{figure}
	
	A figura \ref{fig:tei-gi} tem um rótulo, \texttt{fig:tei-gi}, através do
	qual eu posso referenciá-la em qualquer parte do documento através do
	comando \verb|\ref{fig:tei-gi}|. Se quiser, também posso fazer uma
	referência à página onde esta figura se encontra através do comando
	\verb|\pageref{fig:tei-gi}|. Por exemplo, esta figura está na página
	\pageref{fig:tei-gi}.

	Note também que, ao compilar o arquivo pela primeira vez ou após ter
	alterado ou inserido algum rótulo, o \LaTeX\ dá uma mensagem de aviso,
	informando que os rótulos foram alterados e que você deve re-compilar:

	\begin{center}\slshape
		Label(s) may have changed. Rerun to get cross-references right.
	\end{center}
	
	% Dica: o pacote caption dá várias opções de personalização da legenda
	% de figuras e de tabelas.
	
	\noindent{\color{red}\bfseries Atenção: não coloque o ambiente figure
	dentro de uma caixa!}
	
\end{document}